\documentclass[solutions,answers]{exam}

\usepackage{amsmath}
\usepackage{amssymb}
\usepackage{tabularx}
\usepackage{graphicx}
\usepackage{enumerate}
\usepackage{hhline}
\usepackage{tikz}


%######################################
%Various shortcuts
%Some may need package amssymb

%Shortcuts for various sets
\providecommand{\A}{\ensuremath{\mathfrak{A}}}%ideal https://www.overleaf.com/project/5ba957d84dc5ce49f189bc04A
\providecommand{\B}{\ensuremath{\mathcal{B}}}%basis
\providecommand{\C}{\ensuremath{\mathbb{C}}}%complex
\providecommand{\F}{\ensuremath{\mathbb{F}}}%finite field
\providecommand{\N}{\ensuremath{\mathbb{N}}}%natural
\providecommand{\RI}{\ensuremath{\mathcal{O}}}%ring of integers
\providecommand{\p}{\ensuremath{\mathfrak{P}}}%prime ideal
\providecommand{\Q}{\ensuremath{\mathbb{Q}}}%rationals
\providecommand{\R}{\ensuremath{\mathbb{R}}}%reals
\providecommand{\T}{\ensuremath{\mathcal{T}}}%topology
\providecommand{\Z}{\ensuremath{\mathbb{Z}}}%integers
\providecommand{\Zpos}{\ensuremath{\mathbb{Z}^{+}}}%positive integers


%Miscellaneous shortcuts
\providecommand{\lcm}{\ensuremath{\text{lcm}}}%least common multiple
\providecommand{\st}{\ensuremath{\backepsilon}}
\providecommand{\ud}{\ensuremath{\,\,\mathrm{d}}}%For derivatives
%######################################

%######################################
% Various operators
% some may need package amsmath

%Algebra
\providecommand{\amod}{\ensuremath{\diagup}}%Algebraic mod
\providecommand{\embed}{\ensuremath{\hookrightarrow}}%inclusion map
\providecommand{\gen}[1]{\ensuremath{\left<#1\right>}}%Group generated by #1
\providecommand{\idx}[2]{\ensuremath{\left[#1:#2\right]}}%index of #2 in #1
\providecommand{\im}{\ensuremath{\mathrm{im}}}%Image
\providecommand{\iso}{\ensuremath{\simeq}}%isomorphism
\providecommand{\normal}{\ensuremath{\vartriangleleft}}%Normal subgroup
\providecommand{\subgp}{\ensuremath{\le}}%Sub group
\providecommand{\vect}[1]{\ensuremath{\left\langle#1\right\rangle}}%Vector

%Logic
\providecommand{\land}{\ensuremath{\wedge}}
\DeclareMathSymbol\lnot{\mathbin}{symbols}{"3A}%p528
\providecommand{\lor}{\ensuremath{\vee}}
\providecommand{\lxor}{\ensuremath{\oplus}}

%Number Theory
\providecommand{\jacobi}[2]{\ensuremath{\left(\frac{#1}{#2}\right)}}
\providecommand{\legendre}[2]{\ensuremath{\left(\frac{#1}{#2}\right)}}

%Set Theory
\providecommand{\intersect}{\ensuremath{\bigcap}}
\providecommand{\less}{\ensuremath{\diagdown}}
\providecommand{\union}{\ensuremath{\bigcup}}

%Miscellaneous
\providecommand{\abs}[1]{\ensuremath{\left\lvert#1\right\rvert}}
\providecommand{\card}[1]{\ensuremath{\left\lvert#1\right\rvert}}
\providecommand{\ceiling}[1]{\ensuremath{\left\lceil#1\right\rceil}}
\providecommand{\define}{\ensuremath{\stackrel{\text{\tiny def}}{=}}}
\providecommand{\floor}[1]{\ensuremath{\left\lfloor#1\right\rfloor}}
\providecommand{\norm}[1]{\ensuremath{\left\lVert#1\right\rVert}}%Analysis norm
\providecommand{\order}[1]{\ensuremath{\left\lvert#1\right\rvert}}

%Asymptotic
\providecommand{\Oh}{\ensuremath{\mathcal{O}}}
\providecommand{\oh}{\ensuremath{o}}

%Environments




\header{COMP 2300: Examples}{Examples\ifprintanswers{}\fi}{Fall 2025}
\lfoot{}
\cfoot{}
\rfoot{\thepage}



%%%%%%%%%%%%%%%%%%%%%%%%%%%DO NOT CHANGE THE ABOVE
%%%%%%%%%%%%%%%%%%%%%%%%%%%HOMEWORK STARTS HERE
\begin{document}
\vspace*{1em}
\noindent Name: \textbf{**YOUR NAME HERE**} \setlength\answerclearance{0.7ex}

\begin{questions}

% ============ 1. Propositional logic + align
\question[4]
(\textbf{Logic \& \texttt{align}}) Use the shortcuts $\land,\ \lor,\ \lnot,\ \lxor$ to simplify the formula
\[
(p \land (p \to q)) \to q
\]
by a sequence of equivalences. Show each step in an \texttt{align} environment.

\begin{solution}
Recall $p\to q \equiv \lnot p \lor q$ and $A\to B \equiv \lnot A \lor B$.
\begin{align}
(p \land (p \to q)) \to q
&\equiv \ \lnot(p \land (p \to q)) \lor q \\
&\equiv \ \lnot p \lor \lnot(p \to q) \lor q \quad (\text{De Morgan})\\
&\equiv \ \lnot p \lor \lnot(\lnot p \lor q) \lor q \\
&\equiv \ \lnot p \lor (p \land \lnot q) \lor q \quad (\text{De Morgan})\\
&\equiv \ (\lnot p \lor p) \lor (\lnot p \lor \lnot q) \lor q \\
&\equiv \ \text{True} \lor (\lnot p \lor \lnot q) \lor q \\
&\equiv \ \text{True}.
\end{align}
Thus the statement is a tautology.
\end{solution}

% ============ 2. Truth table with macros
\question[4]
(\textbf{Truth table}) Build the truth table for $(p \lxor q) \lor (p \land q)$.

\begin{solution}
\[
\begin{array}{c|c|c|c|c}
p & q & p\lxor q & p\land q & (p \lxor q)\lor(p\land q)\\
\hhline{=====}
0 & 0 & 0 & 0 & 0\\
0 & 1 & 1 & 0 & 1\\
1 & 0 & 1 & 0 & 1\\
1 & 1 & 0 & 1 & 1\\
\end{array}
\]
The disjunction is true exactly when at least one of $p,q$ is true, so it is equivalent to $p\lor q$.
\end{solution}

% ============ 3. De Morgan + sets
\question[4]
(\textbf{Sets \& complements}) For sets $A,B$ in a universe $U$, show $\overline{A\union B}=\overline{A}\intersect\overline{B}$ and $\overline{A\intersect B}=\overline{A}\union\overline{B}$, using membership reasoning.

\begin{solution}
$x\in\overline{A\union B}$ iff $x\notin A\union B$ iff $(x\notin A)\land(x\notin B)$ iff $x\in\overline{A}$ and $x\in\overline{B}$, i.e. $x\in \overline{A}\intersect\overline{B}$.  
Similarly, $x\in\overline{A\intersect B}$ iff $x\notin A\intersect B$ iff $(x\notin A)\lor(x\notin B)$ iff $x\in\overline{A}$ or $x\in\overline{B}$, i.e. $x\in \overline{A}\union\overline{B}$.
\end{solution}

% ============ 4. Set difference and identities
\question[4]
(\textbf{Set difference}) Using $\less$ for set difference, prove $(A\less B)\union(A\intersect B)=A$ and $(A\less B)\intersect B=\varnothing$.

\begin{solution}
By membership: $x\in A\less B$ iff $x\in A$ and $x\notin B$.  
First identity: The LHS collects all $x\in A$ that are either outside $B$ or inside $B$, hence exactly all of $A$.  
Second identity: If $x\in (A\less B)\intersect B$ then $x\in B$ and $x\notin B$, impossible. So the intersection is empty.
\end{solution}

% ============ 5. Quantifiers and negation
\question[4]
(\textbf{Quantifiers}) Negate and simplify: $\forall x\in \Z\, \exists y\in \Z\ (x=2y \lor x=2y+1)$.

\begin{solution}
\begin{align}
\lnot\Big(\forall x\in \Z\, \exists y\in \Z\ (x=2y \lor x=2y+1)\Big)
&\equiv \exists x\in \Z\, \lnot\Big(\exists y\in \Z\ (x=2y \lor x=2y+1)\Big)\\
&\equiv \exists x\in \Z\, \forall y\in \Z\, \lnot(x=2y \lor x=2y+1)\\
&\equiv \exists x\in \Z\, \forall y\in \Z\, \big((x\ne 2y)\land(x\ne 2y+1)\big).
\end{align}
This says “there is an integer not even and not odd,” which is false—so the original statement is true.
\end{solution}

% ============ 6. Functions (domain/codomain, injective/surjective) + align
\question[6]
(\textbf{Functions}) Let $f:\Z\to \Z$ be $f(n)\define 2n+1$.
\begin{parts}
\part Show $f$ is injective.
\begin{solution}
If $f(a)=f(b)$ then $2a+1=2b+1$, hence $2a=2b$ and $a=b$. So $f$ is injective.
\end{solution}
\part Show $f$ is not surjective onto $\Z$.
\begin{solution}
All values of $f$ are odd. No even integer is in $\im(f)$. Therefore $f$ is not surjective onto $\Z$.
\end{solution}
\part Compute $(f\circ f)(n)$ using \texttt{align}.
\begin{solution}
\begin{align}
(f\circ f)(n) &= f\big(f(n)\big) = f(2n+1) = 2(2n+1)+1 \\
&= 4n+3.
\end{align}
\end{solution}
\end{parts}

% ============ 7. Floors, ceilings, absolute value, norms
\question[4]
(\textbf{Floor/Ceiling/Abs/Norm}) Give numerical examples:
\begin{parts}
\part Compute $\floor{3.9}$, $\ceiling{3.1}$, $\abs{-7}$.
\begin{solution}
$\floor{3.9}=3$, $\ceiling{3.1}=4$, $\abs{-7}=7$.
\end{solution}
\part If $\vec v=(3,-4)$, compute $\norm{\vec v}$.
\begin{solution}
$\norm{\vec v}=\sqrt{3^2+(-4)^2}=\sqrt{25}=5$.
\end{solution}
\end{parts}

% ============ 8. Modular arithmetic (a standard discrete tool)
\question[4]
(\textbf{Modular arithmetic}) Prove: if $a\equiv b\pmod{m}$ and $c\equiv d\pmod{m}$, then $a+c\equiv b+d\pmod{m}$.

\begin{solution}
$a\equiv b\pmod m$ means $m\mid(a-b)$; $c\equiv d\pmod m$ means $m\mid(c-d)$. Then
\[
(a+c)-(b+d)=(a-b)+(c-d),
\]
a sum of multiples of $m$, hence divisible by $m$. Therefore $a+c\equiv b+d\pmod m$.
\end{solution}

% ============ 9. Induction with align
\question[6]
(\textbf{Induction \& \texttt{align}}) Prove by induction that for $n\in\N$,
\[
\sum_{k=1}^n k = \frac{n(n+1)}{2}.
\]

\begin{solution}
\textit{Base $n=1$:} $\sum_{k=1}^1 k=1=\frac{1\cdot2}{2}$.  
\textit{Step:} Assume $\sum_{k=1}^n k=\frac{n(n+1)}{2}$. Then
\begin{align}
\sum_{k=1}^{n+1} k &= \left(\sum_{k=1}^{n} k\right) + (n+1)
= \frac{n(n+1)}{2} + (n+1) \\
&= \frac{n(n+1)+2(n+1)}{2}
= \frac{(n+1)(n+2)}{2}.
\end{align}
Thus the formula holds for all $n\in\N$.
\end{solution}

% ============ 10. Relations (equivalence relation)
\question[6]
(\textbf{Relations}) Define $x\sim y$ on $\Z$ iff $x\equiv y\pmod 3$.
\begin{parts}
\part Prove $\sim$ is an equivalence relation.
\begin{solution}
Reflexive: $x\equiv x\pmod 3$. Symmetric: if $x\equiv y$ then $y\equiv x$. Transitive: if $x\equiv y$ and $y\equiv z$, then $x\equiv z$ (add the congruences). Hence $\sim$ is an equivalence relation.
\end{solution}
\part List the three equivalence classes.
\begin{solution}
$[0]=\{\,\ldots,-6,-3,0,3,6,\ldots\}$, $[1]=\{\,\ldots,-5,-2,1,4,7,\ldots\}$, $[2]=\{\,\ldots,-4,-1,2,5,8,\ldots\}$.
\end{solution}
\end{parts}

% ============ 11. Graph notation (standard in discrete math)
\question[4]
(\textbf{Graphs}) Let $G=(V,E)$ be a simple graph with $V=\{1,2,3,4\}$ and $E=\{\{1,2\},\{2,3\},\{3,4\},\{4,1\}\}$. Compute the degree of each vertex and verify the Handshaking Lemma.

\begin{solution}
Each vertex has degree $2$. Thus $\sum_{v\in V}\deg(v)=8=2|E|=2\cdot4$. Handshaking Lemma holds.
\end{solution}

% ============ 12. Asymptotics example
\question[4]
(\textbf{Asymptotics}) Show $3n^2+10n+7\in \Oh(n^2)$ by finding $C>0$ and $n_0$ so that for all $n\ge n_0$,
\[
3n^2+10n+7 \le C n^2.
\]

\begin{solution}
For $n\ge 1$, $10n\le 10n^2$ and $7\le 7n^2$. Hence
\[
3n^2+10n+7 \le 3n^2+10n^2+7n^2 = 20n^2.
\]
Take $C=20$ and $n_0=1$. Therefore $3n^2+10n+7\in \Oh(n^2)$.
\end{solution}

\end{questions}
\end{document}